\documentclass[12pt, unicode, a4paper]{ltjsreport}

\usepackage[top=25truemm, bottom=25truemm, left=20truemm, right=20truemm]{geometry}
\usepackage[unicode, hidelinks, bookmarksnumbered, pdfusetitle]{hyperref}
\usepackage{chngcntr}
\usepackage{tocloft}

\counterwithout{section}{chapter}
\renewcommand{\thesection}{第\arabic{section}条}

\cftsetindents{chapter}{0truemm}{20truemm}
\cftsetindents{section}{5truemm}{20truemm}

\title{名工大写真映像部 部内規約（部則）}
\date{
  \begin{tabular}{rl}
    制定： & 2026年4月1日 \\
    最終改定： & 2026年9月1日
  \end{tabular}
}
\begin{document}

\maketitle

\tableofcontents

\chapter{総則}
    \section{名称}
        \begin{enumerate}
            \item 当部は名古屋工業大学公認課外活動団体であり、「名工大写真映像部（nitechphoto）」と称す。
        \end{enumerate}

    \section{設立日}
        \begin{enumerate}
            \item 当部の設立日は2026年4月1日である。
        \end{enumerate}

    \section{所在地}
        \begin{enumerate}
            \item 当部の所在地は 〒466-8555 名古屋市昭和区御器所町名古屋工業大学55号館2階暗室 である。ただし、郵便その他配達物については 同大学19号館学生センター横の郵便ポスト にて受け付ける。
        \end{enumerate}

    \section{目的}\label{purpose}
        \begin{enumerate}
            \item 当部は、写真・映像に関する技術の向上と知識の研鑽に努め、交流することを目的とする。
        \end{enumerate}

    \section{存続}
        \begin{enumerate}
            \item 当部は、部員または役員の加入、退部、交代、その他の変更にかかわらず、 同一の団体として存続する。
        \end{enumerate}

\chapter{入部・退部}
    \section{入部}
        \begin{enumerate}
            \item 当部は以下のすべての要件を満たした者による入部を受け付ける。
            \begin{itemize}
                \item 名古屋工業大学または名古屋工業大学大学院に在学中である。
                \item 当規約\ref{budget}2項によって議決された金額を部費として支払った。
                \item 以下のすべての情報を部に提出した。
                \begin{itemize}
                    \item 氏名
                    \item 学籍番号
                    \item 現住所
                    \item メールアドレス
                \end{itemize}
                \item 当規約に同意した。
                \item 当部の定める連絡サービス等のグループに参加している。
            \end{itemize}
        \end{enumerate}

    \section{退部}
        \begin{enumerate}
            \item 役員を除くすべての部員は代表または役員に退部の意思を表明することで任意の時期に退部することができる。
            \item 役員は当規約\ref{resignation}の規定に基づき役員を辞任した場合、前項の規定に基づき退部することができる。
            \item 部員が死亡したときは退部したものとみなす。
        \end{enumerate}


\chapter{懲戒}
    \section{懲戒}\label{disciplinary}
        \begin{enumerate}
            \item 全役員が出席する部会においてその出席者の3分の2以上の承認が得られた場合、当部は部員に対して以下の懲戒処分を下すことができる。ただし、部員資格剥奪（強制退部）を行う際は役員全員の承認を必要とする。
            \begin{itemize}
                \item 厳重注意：代表より厳重に注意を行う。
                \item 戒告：部会にて顛末を説明し、役員より厳重注意を行う。
                \item 部員資格停止：部員としての資格を一時的に停止する。停止期間中は部に関わる設備を利用できない。
                \item 部員資格剥奪（強制退部）：懲戒対象の意志に関わらず部員の資格を喪失する。
            \end{itemize}
            \item 懲戒事由は以下の通りとする。
            \begin{itemize}
                \item 公序良俗・法令等に違反し、反社会的な行動を行った。
                \item 正当な理由なく部費を支払わなかった。
                \item 上記のほか、部の秩序を著しく乱すようなことを行った。
            \end{itemize}
            \item 懲戒審議に際して、懲戒対象となり得る部員は部会または役員会にて釈明を行うことができる。釈明に対して、部会または役員会の出席者は十分な配慮と酌量を行わなければならない。
            \item 懲戒を行うことが議決されてから実際に懲戒を行うまでには、1週間以上の猶予期間を設けなければならない。
            \item 代表は懲戒対象となった部員に対し、懲戒を行う旨を猶予期間中に合理的な手段で通達しなければならない。
            \item 通達なしでの懲戒は無効とする。ただし、懲戒対象の部員が合理的期間内に通達への応答がない場合、これは通達したものとみなす。
        \end{enumerate}


\chapter{役員}
    \section{役員・役職}\label{directors}
        \begin{enumerate}
            \item 当部は部の円滑な運営のために役員をおくものとする。各役員が務める役職とその定員は以下の通りとする。
            \begin{itemize}
                \item 代表（1名）：当部の代表
                \item 副代表（1名程度）：代表の補佐
                \item 会計（1名程度）：部費の管理
                \item 庶務（2名程度）：行事の運営
            \end{itemize}
            \item 役員の任期は当規約\ref{activity}1項にて定められた活動期間とする。
            \item 部会にて当規約\ref{decision}に基づいた承認を得ることで、新たな役職をおくこと、役員の定員を変更することができる。
            \item 各役員は他の役職を兼任することができる。
        \end{enumerate}

    \section{役員の選任}\label{election}
        \begin{enumerate}
            \item 次年の役員の選任は、部員の互選によって選出し、当規約\ref{decision}に基づいた承認により選任する。
            \item 前項にて規定された互選によって選出しない場合は、前年度の代表または役員会が候補者を推薦し、当規約\ref{decision}に基づいた承認を得る。
            \item 前項の規定により推薦を受けた部員はやむを得ない状況を除いて推薦を受諾する。
            \item 役員の再選は妨げないものとする。
        \end{enumerate}

    \section{役員の辞任・解任}\label{resignation}
        \begin{enumerate}
            \item 代表は、役員会で辞任の意思を表明し、全役員が出席する部会においてその出席者の過半数の承認が得られた場合、代表を辞任することができる。
            \item 代表以外の役員は、役員会で辞任の意思を表明し、部会にて当規約\ref{decision}に基づいた承認が得られた場合、辞任することができる。
            \item 当規約\ref{disciplinary}に違反するなど役員として不適格と認められた者は、当規約\ref{decision}に基づいた承認が得られた場合、強制的に解任させられる。
            \item 前項の規定により役員が解任される場合、後任の役員は時期に関わらず、直ちに当規約\ref{election}に基づき選任される。
        \end{enumerate}

    \section{役員会}\label{directors-meeting}
        \begin{enumerate}
            \item 役員会は任意の役員が必要に応じ他の役員を召集することによって不定期に開催される。
            \item 役員会における承認・議決は、役員の懲戒・解任審議を除いて、過半数の賛成によって可決する。
            \item 必要に応じて、役員以外の部員、OB/OG、その他を役員会に参加させることができる。ただし、その者の議決権は原則ないものとする。
        \end{enumerate}


\chapter{OB/OG}
    \section{OB/OG}
        \begin{enumerate}
            \item 名古屋工業大学または名古屋工業大学大学院を卒業または学籍を離れ、在学中に当部に部員として所属していた者をOB/OGとする。
        \end{enumerate}


\chapter{活動}
    \section{活動}\label{activity}
        \begin{enumerate}
            \item 当部の活動期間は毎年3月1日より2月末までとする。
            \item 当部は当規約\ref{purpose}の目的を達するために以下の活動を行う。
            \begin{itemize}
                \item 撮影会
                \item 展示会
                \item 当規約\ref{meeting}によって規定される部会
                \item 当規約\ref{event}によって規定される行事
                \item 新入生の勧誘活動
            \end{itemize}
            \item 全ての部員は当規約\ref{purpose}の目的を達するため、前項の活動に積極的に参加することが求められるが、学業を優先させることは妨げられない。
        \end{enumerate}

    \section{行事}\label{event}
        \begin{enumerate}
            \item 当部は行事として以下のものを開催する。
            \begin{itemize}
                \item 新歓BBQ
                \item 夏合宿
                \item 忘年会
                \item その他のイベント
            \end{itemize}
        \end{enumerate}

    \section{制作物の権利・利用}\label{production}
        \begin{enumerate}
            \item 当部の活動に関連して部員が制作した写真、映像、その他の制作物（以下「制作物」という）の著作権は、原則として当該制作物を制作した部員に帰属するものとする。
            \item 制作物が複数の部員による共同制作である場合は、当該制作物の著作権は当該部員に共有されるものとする。
            \item 前項の規定にかかわらず、当規約\ref{purpose}の目的を達するために必要な範囲（部の運営、広報、記録、展示、勧誘活動、その他）において、
            当部は当該制作物を無償で利用（複製、展示、公衆送信、SNSへの掲載を含む）することができるものとし、制作した部員は、これをあらかじめ承諾するものとする。
            制作物の利用にあたっては、著作者人格権に十分配慮するものとする。
            \item 当部が前項の目的を超えて制作物を利用する場合、または第三者に利用させる場合には、原則として当該制作物を制作した部員の事前の同意を得なければならない。
            \item 制作物の利用にあたっては、特段の事情がない限り、制作者の氏名またはハンドルネーム等を適切に表示するよう努めるものとする。
            \item 部員が退部した後であっても、在部中に制作された制作物については、本条の規定が引き続き適用されるものとする。
            ただし、退部後に使用の取り止めを求められた際は、当部はその意思を尊重し、合理的な範囲（当部の運営に著しい支障を及ぼさない範囲など）で対応するものとする。
            \item 制作物に第三者の肖像が含まれる場合、当部および部員は肖像権、その他の権利に十分配慮し、必要に応じて被写体の同意を得るものとする。
        \end{enumerate}


\chapter{部会}
    \section{部会}\label{meeting}
        \begin{enumerate}
            \item 部会は全ての部員をもって構成し、当部の最高意思決定機関とする。部会は、一般にいう総会に相当するものとする。
            \item 部会は以下のことを行うために開催される。
            \begin{itemize}
                \item 当部の運営のために必要な審議・議決
                \item 役員の選任および解任
                \item 会計の承認
                \item 規約の制定および改定
                \item 活動報告
                \item 連絡事項の通達
            \end{itemize}
            \item 部会は原則毎月決まった曜日に開催されるものとする。なお、開始時刻・開催場所は部会ごとに役員が決定する。
            \item 役員は必要に応じて臨時部会を開催することができる。
        \end{enumerate}

    \section{部会における審議・議決}\label{decision}
        \begin{enumerate}
            \item 部会の議決は他に定めのない限り、部員の3分の1以上が出席する部会において、その出席者の過半数の賛成によって可決される。
            \item 懲戒審議と役員の解任審議は他のどのような審議よりも優先される。
            \item 部会における審議の際の発言の責任は、著しく部の秩序を乱すもので無い限り問われないものとする。
            \item 部会における議決の際は、当規約とは別に定める書式に基づく議決書を原則作成するものとする。
            \item 当規約における「正当な理由」の有無は、原則として代表または役員会がこれを判断する。
        \end{enumerate}

    \section{部会のみなし出席・委任等}\label{meeting-decision}
        \begin{enumerate}
            \item 当規約において「全役員が出席する部会」は、全役員が現に出席している場合のほか、
            正当な理由により欠席する役員があらかじめ電磁的手段または書面によって、委任または当該議題に対する意思表示（賛否を含む）のいずれかを行った場合には、
            その役員を出席したものとみなし、当該部会における出席者数および議決権の算定に含めるものとして、当該意思表示または委任の範囲において議決権を有するものとする。
            \item ただし、やむを得ない事情により委任または意思表示を得ることができない役員がある場合は、役員会の判断により当該役員を除いた全役員の出席をもって足りるものとする。
            \item 委任は、当該部会における議題ごと、または包括的に行うことができる。また委任および意思表示の双方が行われた場合は、当該議題に対する意思表示を優先するものとする。
        \end{enumerate}


\chapter{部費}
    \section{部費}\label{budget}
        \begin{enumerate}
            \item 部費は以下の物品を購入するために使用される。
            \begin{itemize}
                \item 写真用品
                \item 機材
                \item 当部が活動する上で必要となる日用品
                \item その他、必要となる物品・サービス等
            \end{itemize}
            \item 部費の金額は当規約\ref{activity}1項にて定められた活動期間の最初の月の部会までに役員会において案を作成して、\ref{decision}に基づいた承認を得る。
        \end{enumerate}

    \section{部費の管理・監査}
        \begin{enumerate}
            \item 部費の管理は当規約\ref{directors}にて規定される会計が行う。
            \item 部費の管理を円滑に行うために当部は以下の銀行口座を保有する。
            \begin{itemize}
                \item 設立後に速やかに開設して周知する。
                % \item xxx銀行 xxxx
            \end{itemize}
            \item 会計は当規約\ref{activity}1項にて定められた活動期間の最後の月の部会にて会計報告を行い、\ref{decision}に基づいた承認を得る。
            \item 部費の監査は会計以外の役員が行う。
            \item 部費に損害が生じた場合は、原則として役員の故意または重過失による場合に限り責任を負う。
        \end{enumerate}

    \section{部費の徴収}
        \begin{enumerate}
            \item 新入部員以外の部員は当規約\ref{budget}2項によって議決された金額を、部費として当規約\ref{activity}1項にて規定された活動期間の最初の月中に当部の銀行口座へ振り込まなければならない。
            \item 新入部員は当規約\ref{budget}2項によって議決された金額を、部費として入部後1ヶ月以内に当部の銀行口座へ振り込まなければならない。
            \item 追加の部費を徴収する場合は、全役員が出席する部会において、会計がその時点までの会計報告を行った上で、その出席者の過半数の承認を得なければならない。
            \item 残金は次年の部費に積み立てる。
            \item 退部時の部費の返却は一切行わないものとする。ただし、部会にて当規約\ref{decision}の規定に基づいた承認が得られた場合はその限りではない。
        \end{enumerate}

    \section{部費による物品やサービスの購入}
        \begin{enumerate}
            \item 部費によって物品やサービスを購入する際は、事前または事後に役員会における承認を必要とする。
            \item 部費によって物品やサービスを購入する際に立替が必要となる場合、その立替は任意の部員が行う。
            \item 会計は前項の規定に基づき立替を行った部員に対し、可能な限り早く無利子で返金するものとする。
        \end{enumerate}


\chapter{改定}
    \section{改定}
        \begin{enumerate}
            \item 当規約は全役員が出席する部会においてその出席者の過半数の承認が得られた場合、改定することができる。ただし、当該部会を開催する1週間以上前に連絡サービス等で改定箇所を示さなければならない。
            \item 改定規則は原則即時施行とする。ただし、部会の議決によって公布期間を設けることができる。
        \end{enumerate}


\chapter{免責}
    \section{免責}
        \begin{enumerate}
            \item 当部は重過失が認められない限り、当部の活動に伴い生じた事故や事件等に関して法令に反しない範囲でいかなる責任も負わず、各個人の責任とする。ただし、部会にて当規約\ref{decision}に基づいた承認が得られた場合はこの限りではない。
        \end{enumerate}
    \section{保険}
        \begin{enumerate}
            \item 部員が当規約\ref{activity}に規定される活動を行うときは不慮の事故等に備え、損害賠償保険や人身傷害保険などの大学指定の保険または当部が必要と認める保険に加入しなければならない。
            \item 保険に加入しない部員に対して役員はその部員の活動を制限することができる。
        \end{enumerate}

    \section{個人情報・データの取扱い}
        \begin{enumerate}
            \item 当部の個人情報の取り扱いについては、原則「個人情報の保護に関する法律」第1条の目的を達することを目標とし、厳重にこれを管理する。
            取得した個人情報は、当部の運営および活動に必要な範囲でのみ利用するものとする。
            \item 書面または電磁的記録媒体で保存されているデータは可能な範囲でバックアップを行う。
            \item 部員は電磁的記録媒体で保存されているデータが喪失したとしても、故意または重過失がない限り責任を負わない。
            \item 部員が退部した場合は退部した部員に関わる個人情報・データは原則速やかに破棄される。ただし、当該部員との合意が得られた場合はこの限りではない。
        \end{enumerate}


\chapter{その他}
    \section{その他}
        当規約に定めるほかに必要な事項は十分に配慮の上、役員会・部会でこれを別に定める。

    \section*{制定}
        当規約は、制定日である 2026年4月1日から発効する。
        \begin{itemize}
            \item 2026年4月1日：当規約の制定
            \item 2026年9月1日：当規約の大幅な改定
        \end{itemize}

    \section*{代表者による証明}
        当規約の記載内容について、事実と相違ないことをここに証明する。
        \begin{itemize}
            \item 代表者
            \begin{flushright}
                （印）
            \end{flushright}
            \item 住所
        \end{itemize}

\end{document}
